%*******************************************************
% Abstract
%*******************************************************
%\renewcommand{\abstractname}{Abstract}
\pdfbookmark[1]{Abstract}{Abstract}
% \addcontentsline{toc}{chapter}{\tocEntry{Abstract}}
\begingroup
\let\clearpage\relax
\let\cleardoublepage\relax
\let\cleardoublepage\relax

\chapter*{Abstract}
% Short summary of the contents in English\dots a great guide by
% Kent Beck how to write good abstracts can be found here:
% \begin{center}
% \url{https://plg.uwaterloo.ca/~migod/research/beckOOPSLA.html}
% \end{center}

% THIS IS ONLY A DRAFT
Large Language Models encode their computations in dense, high-dimensional
activations whose individual dimensions seldom correspond to anything a human
can name, and Sparse AutoEncoders (SAEs) have become the standard tool for
decomposing these activations into interpretable concepts. Yet SAEs impose a
steep sparsity–reconstruction trade-off: because each concept is a single scalar
coefficient, holding the reconstruction error at an acceptable level leaves far
too many active concepts per token for the decomposition to be interpretable.
This thesis introduces the Variational Auto Embedding Encoder (VAEE), a model
that represents each concept not as a scalar but as a learned vector embedding,
drawing together ideas from Concept Embedding Models, Variational AutoEncoders,
and Vector-Quantized Variational AutoEncoders (VQ-VAEs); like a VQ-VAE it
decodes through a learned codebook, but rather than reconstructing from a single
embedding it stacks every embedding deemed active --- each selected by comparing
the encoded input against a learned prototype --- and decodes from that stack.
Derived from first principles, the VAEE markedly relaxes the
sparsity–reconstruction trade-off: on Mistral 7B residual-stream activations, a
VAEE with eight active concepts matches the reconstruction of a
parameter-matched TopK SAE that requires 128. These gains come with honest
caveats: reconstruction was assessed only by mean squared error, each concept
now spans several active neurons rather than one, and the recovered concepts
were not found to be more interpretable than the SAE's --- so the VAEE's
contribution lies in expanding the capacity and sparsity of concept
representations, leaving their interpretability and downstream steering as the
central open questions.
% END DRAFT


% \vfill

% \begin{otherlanguage}{ngerman}
% \pdfbookmark[1]{Zusammenfassung}{Zusammenfassung}
% \chapter*{Zusammenfassung}
% Kurze Zusammenfassung des Inhaltes in deutscher Sprache\dots
% \end{otherlanguage}

\endgroup

\vfill
